% ===========================================================================
% SOURCES: docs/claude_project/11_REPO_MAP.md
%          scripts/README.md, firmware/README.md, README.md
%
% The other half of what "Appendix B: Source Code" was trying to do, done
% usefully: not the code, but enough to re-run it.
% ===========================================================================
\chapter*{Appendix C: Building and Reproducing the Results}
\label{appendix:reproduction}
\addcontentsline{toc}{chapter}{Appendix C: Building and Reproducing the Results}

\section{Repository Layout}
\label{sec:repo-layout}
% One table, one line per top-level directory. From 11_REPO_MAP.md.

\section{Reproducing the Software Model}
\label{sec:repro-sw}
% Training, quantisation, export. State the checkpoint hash (\thCkptSha) and
% that it is tracked in Git LFS - a reader who gets a different hash is not
% reproducing this thesis and should know immediately. (After DD-036's
% redeploy that hash is the leaf-split checkpoint's, \thAccCkptSha; check that
% the two macros agree before quoting either.)

\section{Reproducing the Accuracy Evaluation}
\label{sec:repro-accuracy}
% \DD{034}. software/train/resplit_dataset.py rebuilds data/grape_leafsplit
% from the downloaded images and PlantVillage's leaf maps - not committed, so
% say where to fetch them and which upstream commit they are pinned to; then
% train.py, quantize.py, and `evaluate.py --split test --report', whose JSON
% is what thesis_numbers.py reads. The test split is scored once: say so, so
% nobody "reproduces" it by tuning on it.

\section{Running the Simulations}
\label{sec:repro-sim}
% scripts/run_sim.sh for one testbench, run_all_sims.sh for all
% \thTestbenches. Say what a pass looks like (ALL PASS and the element count),
% so the claim in Section~\ref{sec:bit-exactness-result} is checkable rather than
% quotable.

\section{Synthesis and Implementation}
\label{sec:repro-impl}
% run_synth.sh and run_impl.sh, with TOP and PERIOD. Note that every number in
% Chapter~\ref{chap:timing} and Chapter~\ref{chap:evaluation} came from these two scripts, and
% that the outputs land in build/<run>/ with the reports the numbers were
% parsed from.

\section{The Block Design and the Bitstream}
\label{sec:repro-bd}
% run_bd.sh, make_bd.tcl, and where the XSA and the bitstream land.

\section{Building and Staging the PS Application}
\label{sec:repro-ps}
% firmware/Makefile: the host unit tests, `make lint', and the bare-metal
% build. Then make_sdcard.py and the known-answer file it stages.

\section{The Live Camera Image}
\label{sec:repro-camera}
% \DD{035}. firmware/linux/petalinux/build.sh on a Linux host (PetaLinux does
% not run on Windows), the static vinecam binary, and the first board step:
% vinecam with a known-answer image and its EXPECTED.txt, before any camera.
% firmware/linux/README.md has the commands; point at it rather than copy it.

\section{The Edge-Device Measurements}
\label{sec:repro-edge}
% \DD{036}. The ONNX export (opset 13) and its check against PyTorch on the
% test images; ONNX Runtime on the Raspberry Pi 4; trtexec on the Jetson Nano
% (the engine is built on the device); the C model built natively on the Pi;
% the power-meter procedure; and the JSON reports thesis_numbers.py reads.
% Give the software versions and the clock settings - a different governor or
% power mode is a different measurement.

\section{Regenerating the Numbers in This Document}
\label{sec:repro-numbers}
% scripts/thesis_numbers.py -> thesis/generated/numbers.tex, and the rule that
% every figure in the manuscript comes from there or from manual_numbers.tex.
% Worth one honest paragraph on why: a number copied by hand is true only until
% the thing it measured changes, which is \DD{019} restated for prose.
